\documentclass[11pt]{article}

\usepackage{fullpage}
\usepackage{amsmath, amssymb, bm, cite, epsfig, psfrag}
\usepackage{graphicx}
\usepackage{float}
\usepackage{amsthm}
\usepackage{amsfonts}
\usepackage{listings}
\usepackage{cite}
\usepackage{hyperref}
\usepackage{tikz}
\usepackage{enumerate}
\usepackage{listings}
\usepackage{mathtools}
\usepackage{mdframed}
\lstloadlanguages{Python}
\usetikzlibrary{shapes,arrows}
%\usetikzlibrary{dsp,chains}

\DeclareFixedFont{\ttb}{T1}{txtt}{bx}{n}{9} % for bold
\DeclareFixedFont{\ttm}{T1}{txtt}{m}{n}{9}  % for normal
% Defining colors
\usepackage{color}
\definecolor{deepblue}{rgb}{0,0,0.5}
\definecolor{deepred}{rgb}{0.6,0,0}
\definecolor{deepgreen}{rgb}{0,0.5,0}
\definecolor{backcolour}{rgb}{0.95,0.95,0.92}

%\restylefloat{figure}
%\theoremstyle{plain}      \newtheorem{theorem}{Theorem}
%\theoremstyle{definition} \newtheorem{definition}{Definition}

\def\del{\partial}
\def\ds{\displaystyle}
\def\ts{\textstyle}
\def\beq{\begin{equation}}
\def\eeq{\end{equation}}
\def\beqa{\begin{eqnarray}}
\def\eeqa{\end{eqnarray}}
\def\beqan{\begin{eqnarray*}}
\def\eeqan{\end{eqnarray*}}
\def\nn{\nonumber}
\def\binomial{\mathop{\mathrm{binomial}}}
\def\half{{\ts\frac{1}{2}}}
\def\Half{{\frac{1}{2}}}
\def\N{{\mathbb{N}}}
\def\Z{{\mathbb{Z}}}
\def\Q{{\mathbb{Q}}}
\def\R{{\mathbb{R}}}
\def\C{{\mathbb{C}}}
\def\argmin{\mathop{\mathrm{arg\,min}}}
\def\argmax{\mathop{\mathrm{arg\,max}}}
%\def\span{\mathop{\mathrm{span}}}
\def\diag{\mathop{\mathrm{diag}}}
\def\x{\times}
\def\limn{\lim_{n \rightarrow \infty}}
\def\liminfn{\liminf_{n \rightarrow \infty}}
\def\limsupn{\limsup_{n \rightarrow \infty}}
\def\MID{\,|\,}
\def\MIDD{\,;\,}

\newtheorem{proposition}{Proposition}
\newtheorem{definition}{Definition}
\newtheorem{theorem}{Theorem}
\newtheorem{lemma}{Lemma}
\newtheorem{corollary}{Corollary}
\newtheorem{assumption}{Assumption}
\newtheorem{claim}{Claim}
\def\qed{\mbox{} \hfill $\Box$}
\setlength{\unitlength}{1mm}

\def\bhat{\widehat{b}}
\def\ehat{\widehat{e}}
\def\phat{\widehat{p}}
\def\qhat{\widehat{q}}
\def\rhat{\widehat{r}}
\def\shat{\widehat{s}}
\def\uhat{\widehat{u}}
\def\ubar{\overline{u}}
\def\vhat{\widehat{v}}
\def\xhat{\widehat{x}}
\def\xbar{\overline{x}}
\def\zhat{\widehat{z}}
\def\zbar{\overline{z}}
\def\la{\leftarrow}
\def\ra{\rightarrow}
\def\MSE{\mbox{\small \sffamily MSE}}
\def\SNR{\mbox{\small \sffamily SNR}}
\def\SINR{\mbox{\small \sffamily SINR}}
\def\arr{\rightarrow}
\def\Exp{\mathbb{E}}
\def\var{\mbox{var}}
\def\Tr{\mbox{Tr}}
\def\tm1{t\! - \! 1}
\def\tp1{t\! + \! 1}
\def\Tm1{T\! - \! 1}
\def\Tp1{T\! + \! 1}


\def\Xset{{\cal X}}

\newcommand{\one}{\mathbf{1}}
\newcommand{\abf}{\mathbf{a}}
\newcommand{\bbf}{\mathbf{b}}
\newcommand{\dbf}{\mathbf{d}}
\newcommand{\ebf}{\mathbf{e}}
\newcommand{\gbf}{\mathbf{g}}
\newcommand{\hbf}{\mathbf{h}}
\newcommand{\pbf}{\mathbf{p}}
\newcommand{\pbfhat}{\widehat{\mathbf{p}}}
\newcommand{\qbf}{\mathbf{q}}
\newcommand{\qbfhat}{\widehat{\mathbf{q}}}
\newcommand{\rbf}{\mathbf{r}}
\newcommand{\rbfhat}{\widehat{\mathbf{r}}}
\newcommand{\sbf}{\mathbf{s}}
\newcommand{\sbfhat}{\widehat{\mathbf{s}}}
\newcommand{\ubf}{\mathbf{u}}
\newcommand{\ubfhat}{\widehat{\mathbf{u}}}
\newcommand{\utildebf}{\tilde{\mathbf{u}}}
\newcommand{\vbf}{\mathbf{v}}
\newcommand{\vbfhat}{\widehat{\mathbf{v}}}
\newcommand{\wbf}{\mathbf{w}}
\newcommand{\wbfhat}{\widehat{\mathbf{w}}}
\newcommand{\xbf}{\mathbf{x}}
\newcommand{\xbfhat}{\widehat{\mathbf{x}}}
\newcommand{\xbfbar}{\overline{\mathbf{x}}}
\newcommand{\ybf}{\mathbf{y}}
\newcommand{\zbf}{\mathbf{z}}
\newcommand{\zbfbar}{\overline{\mathbf{z}}}
\newcommand{\zbfhat}{\widehat{\mathbf{z}}}
\newcommand{\Ahat}{\widehat{A}}
\newcommand{\Abf}{\mathbf{A}}
\newcommand{\Bbf}{\mathbf{B}}
\newcommand{\Cbf}{\mathbf{C}}
\newcommand{\Bbfhat}{\widehat{\mathbf{B}}}
\newcommand{\Dbf}{\mathbf{D}}
\newcommand{\Gbf}{\mathbf{G}}
\newcommand{\Hbf}{\mathbf{H}}
\newcommand{\Ibf}{\mathbf{I}}
\newcommand{\Kbf}{\mathbf{K}}
\newcommand{\Pbf}{\mathbf{P}}
\newcommand{\Phat}{\widehat{P}}
\newcommand{\Qbf}{\mathbf{Q}}
\newcommand{\Rbf}{\mathbf{R}}
\newcommand{\Rhat}{\widehat{R}}
\newcommand{\Sbf}{\mathbf{S}}
\newcommand{\Ubf}{\mathbf{U}}
\newcommand{\Vbf}{\mathbf{V}}
\newcommand{\Wbf}{\mathbf{W}}
\newcommand{\Xhat}{\widehat{X}}
\newcommand{\Xbf}{\mathbf{X}}
\newcommand{\Ybf}{\mathbf{Y}}
\newcommand{\Zbf}{\mathbf{Z}}
\newcommand{\Zhat}{\widehat{Z}}
\newcommand{\Zbfhat}{\widehat{\mathbf{Z}}}
\def\alphabf{{\boldsymbol \alpha}}
\def\betabf{{\boldsymbol \beta}}
\def\betabfhat{{\widehat{\bm{\beta}}}}
\def\epsilonbf{{\boldsymbol \epsilon}}
\def\mubf{{\boldsymbol \mu}}
\def\lambdabf{{\boldsymbol \lambda}}
\def\etabf{{\boldsymbol \eta}}
\def\xibf{{\boldsymbol \xi}}
\def\taubf{{\boldsymbol \tau}}
\def\sigmahat{{\widehat{\sigma}}}
\def\thetabf{{\bm{\theta}}}
\def\thetabfhat{{\widehat{\bm{\theta}}}}
\def\thetahat{{\widehat{\theta}}}
\def\mubar{\overline{\mu}}
\def\muavg{\mu}
\def\sigbf{\bm{\sigma}}
\def\etal{\emph{et al.}}
\def\Ggothic{\mathfrak{G}}
\def\Pset{{\mathcal P}}
\newcommand{\bigCond}[2]{\bigl({#1} \!\bigm\vert\! {#2} \bigr)}
\newcommand{\BigCond}[2]{\Bigl({#1} \!\Bigm\vert\! {#2} \Bigr)}
\newcommand{\tran}{^{\text{\sf T}}}
\newcommand{\herm}{^{\text{\sf H}}}
\newcommand{\bkt}[1]{{\langle #1 \rangle}}
\def\Norm{{\mathcal N}}
\newcommand{\vmult}{.}
\newcommand{\vdiv}{./}


% Python style for highlighting
\newcommand\pythonstyle{\lstset{
language=Python,
backgroundcolor=\color{backcolour},
commentstyle=\color{deepgreen},
basicstyle=\ttm,
otherkeywords={self},             % Add keywords here
keywordstyle=\ttb\color{deepblue},
emph={MyClass,__init__},          % Custom highlighting
emphstyle=\ttb\color{deepred},    % Custom highlighting style
stringstyle=\color{deepgreen},
%frame=tb,                         % Any extra options here
showstringspaces=false            %
}}

% Python environment
\lstnewenvironment{python}[1][]
{
\pythonstyle
\lstset{#1}
}
{}

% Python for external files
\newcommand\pythonexternal[2][]{{
\pythonstyle
\lstinputlisting[#1]{#2}}}

% Python for inline
\newcommand\pycode[1]{{\pythonstyle\lstinline!#1!}}

% Solution environment
\definecolor{lightgray}{gray}{0.95}
\newmdenv[linecolor=white,backgroundcolor=lightgray,frametitle=Solution:]{solution}

\begin{document}

\title{Introduction to Machine Learning\\
Problems: Principal Component Analysis (with Solutions)}
\author{Profs. Sundeep Rangan and Yao Wang}
\date{}

\maketitle

\begin{enumerate}

\item Assume that you have 4 samples each with dimension 3, described in the data matrix $X$,
\[
    X = \left[ \begin{array}{ccc}
        3 & 2 & 1 \\
        2 & 4 & 5 \\
        1 & 2 &3 \\
        0 & 2 &5 \\
          \end{array} \right] \quad
 \]
 For the problems below, you may do the calculations in python.  Describe the commands
 you used.
 \begin{enumerate}
 \item Find the sample mean.
 \item Find the sample covariance matrix $Q$.
 \item Find the eigenvalues and eigenvectors. You can use the 
 \pycode{numpy.linalg.eig} to compute eigenvalues and eigenvectors from $Q$.
 \item Find the PCA coefficients corresponding to samples in $X$.
 \item Reconstruct the original samples from the PCA coefficients.
 \item Approximate the samples using principle components corresponding to the two largest eigenvalues.
 \item Verify the sum of reconstruction error squares = sum of squares of skipped PCA coefficients.
 \end{enumerate}

\begin{solution}
\begin{enumerate}
\item The sample mean is computed as:
\begin{python}
import numpy as np
X = np.array([[3, 2, 1], [2, 4, 5], [1, 2, 3], [0, 2, 5]])
mu = np.mean(X, axis=0)
\end{python}
The mean is $\mubf = [1.5, 2.5, 3.5]$.

\item The sample covariance matrix is computed as:
\begin{python}
X_centered = X - mu
Q = np.cov(X_centered.T)
# Or equivalently: Q = (1/(n-1)) * X_centered.T @ X_centered
\end{python}
The covariance matrix is:
\[
Q = \left[ \begin{array}{ccc}
    2.0 & 0.0 & -1.0 \\
    0.0 & 1.0 & 0.0 \\
    -1.0 & 0.0 & 2.0
\end{array} \right]
\]

\item Eigenvalues and eigenvectors:
\begin{python}
eigenvals, eigenvecs = np.linalg.eig(Q)
# Sort by eigenvalues in descending order
idx = eigenvals.argsort()[::-1]
eigenvals = eigenvals[idx]
eigenvecs = eigenvecs[:, idx]
\end{python}
The eigenvalues are approximately $\lambda_1 = 3.0$, $\lambda_2 = 2.0$, $\lambda_3 = 0.0$.
The corresponding eigenvectors (principal components) are the columns of the eigenvector matrix.

\item PCA coefficients are computed by projecting centered data onto the principal components:
\begin{python}
V = eigenvecs  # Principal components (columns)
Z = X_centered @ V  # PCA coefficients
\end{python}
Each row of $Z$ contains the PCA coefficients for the corresponding sample.

\item Reconstruction from PCA coefficients:
\begin{python}
X_reconstructed = Z @ V.T + mu
\end{python}
This should recover the original $X$ exactly (within numerical precision).

\item Approximation using two largest PCs:
\begin{python}
V2 = eigenvecs[:, :2]  # First two principal components
Z2 = X_centered @ V2  # Coefficients for first two PCs
X_approx = Z2 @ V2.T + mu  # Approximation
\end{python}

\item Verification:
\begin{python}
reconstruction_error = X - X_approx
error_squared_sum = np.sum(reconstruction_error**2)

skipped_coeffs = Z[:, 2:]  # Skipped PCA coefficients
skipped_squared_sum = np.sum(skipped_coeffs**2)

# These should be equal (or very close)
\end{python}
The sum of squares of reconstruction errors equals the sum of squares of the skipped PCA coefficients because the reconstruction error is exactly the projection onto the discarded principal components.
\end{enumerate}
\end{solution}
 
\item \emph{PC bases.}  Suppose that the mean and first two PCs in a basis
 are,
\[
    \mubf = [1,0,2], \quad \vbf_1 = \frac{1}{\sqrt{2}}[1,1,0], \quad
    \vbf_2 = \frac{1}{\sqrt{2}}[1,-1,0], 
\]
\begin{enumerate}[(a)]
  \item What are the two PC coefficients of the vector $\xbf=[2,3,4]$?
  \item What is $\xbfhat$, the approximation of $\xbf$ using two PCs?
  \item What is the approximation error $\|\xbf-\xbfhat\|^2$?
\end{enumerate}

\begin{solution}
\begin{enumerate}[(a)]
\item The PC coefficients are computed by projecting the centered vector onto the PCs:
\[
\xbf_{\text{centered}} = \xbf - \mubf = [2,3,4] - [1,0,2] = [1,3,2]
\]
The coefficients are:
\begin{align*}
z_1 &= \vbf_1^T (\xbf - \mubf) = \frac{1}{\sqrt{2}}[1,1,0]^T [1,3,2] = \frac{1}{\sqrt{2}}(1 + 3) = \frac{4}{\sqrt{2}} = 2\sqrt{2} \\
z_2 &= \vbf_2^T (\xbf - \mubf) = \frac{1}{\sqrt{2}}[1,-1,0]^T [1,3,2] = \frac{1}{\sqrt{2}}(1 - 3) = \frac{-2}{\sqrt{2}} = -\sqrt{2}
\end{align*}

\item The approximation is:
\[
\xbfhat = \mubf + z_1 \vbf_1 + z_2 \vbf_2 = [1,0,2] + 2\sqrt{2} \cdot \frac{1}{\sqrt{2}}[1,1,0] + (-\sqrt{2}) \cdot \frac{1}{\sqrt{2}}[1,-1,0]
\]
\[
= [1,0,2] + 2[1,1,0] - [1,-1,0] = [1,0,2] + [2,2,0] - [1,-1,0] = [2,3,2]
\]

\item The approximation error is:
\[
\|\xbf-\xbfhat\|^2 = \|[2,3,4] - [2,3,2]\|^2 = \|[0,0,2]\|^2 = 4
\]
Note that this error equals the square of the third (skipped) PC coefficient, which would be the projection onto the direction perpendicular to the first two PCs.
\end{enumerate}
\end{solution}

\item \emph{Fitting data with PCs}.  You are given python functions
for PCA transform and a binary classifier:
\begin{python}
    mu, V = PCA(X)  # Finds the mean and PCs for a data matrix X
    
    clf = Classifier()
    clf.fit(Z,y)    # Fits a binary classifier from features Z and labels y
    yhat = clf.predict(Z)  # Predicts labels from Z
\end{python}
Given a data matrix \pycode{X} with binary labels \pycode{y},
you wish to build a classifier that uses a PCA transform followed by
the \pycode{Classifier}.  Write python code that uses model validation
to select the optimal number of PCA components to use.  You may assume:
\begin{itemize}
\item You use simple cross validation with one training and test split.
Not You do not need to do $K$-fold validation.
\item You should use roughly 25\% of the samples for test.  Data shuffling
before splitting is not needed.
\end{itemize}

\begin{solution}
\begin{python}
n = len(X)
n_test = int(0.25 * n)
n_train = n - n_test

# Split data
X_train = X[:n_train]
X_test = X[n_test:]
y_train = y[:n_train]
y_test = y[n_test:]

# Find PCA on training data
mu, V = PCA(X_train)

# Try different numbers of components
best_nc = 0
best_acc = 0

for nc in range(1, min(n_train, X.shape[1]) + 1):
    # Transform training data
    X_train_centered = X_train - mu
    Z_train = X_train_centered @ V[:, :nc]
    
    # Fit classifier
    clf = Classifier()
    clf.fit(Z_train, y_train)
    
    # Transform test data
    X_test_centered = X_test - mu
    Z_test = X_test_centered @ V[:, :nc]
    
    # Predict and evaluate
    yhat_test = clf.predict(Z_test)
    acc = np.mean(yhat_test == y_test)
    
    if acc > best_acc:
        best_acc = acc
        best_nc = nc

print(f"Best number of components: {best_nc}")
\end{python}
\end{solution}

\item \emph{PC with images}.  A dataset has 1000 images of size 28 $\times$ 28,
represented as python array \pycode{X} with shape \pycode{(1000,28,28)}.
You are given the following python functions:
\begin{python}
    Y = reshape(X, shape)         # Reshapes X to a shape 
    pca = PCA(n_components = nc)  # Creates a PCA object
    pca.fit(Y)   # Finds the mean and PCs from training data Y
    Z = pca.transform(Y)   # Find coefficients in the PC basis 
    Yhat = pca.inverse_transform(Z)  # Invert the PC transform
\end{python}
Write a few lines of python code to (i) fit a PC model from the first 500 images
with 5 components; 
and (ii) create an array of approximations of the remaining 500 images.

\begin{solution}
\begin{python}
# (i) Fit PCA model from first 500 images
X_train = X[:500]
X_train_flat = reshape(X_train, (500, 28*28))
pca = PCA(n_components=5)
pca.fit(X_train_flat)

# (ii) Create approximations of remaining 500 images
X_test = X[500:]
X_test_flat = reshape(X_test, (500, 28*28))
Z_test = pca.transform(X_test_flat)
X_test_approx_flat = pca.inverse_transform(Z_test)
X_test_approx = reshape(X_test_approx_flat, (500, 28, 28))
\end{python}
\end{solution}
 
\item \emph{PCs using SVDs.}  You are given a python function:
\begin{python}
    U,s, Vtr = svd(Z, full_matrices = False)
\end{python}
which computes an ``economy" SVD.  Given a data matrix \pycode{X},
write python code that:
\begin{enumerate}[(i)]
\item Finds the PCs and mean of the data.
\item Finds the minimum number of PCs in order that the proportion of variance is
at least 90\%.
\item Create an approximation \pycode{Xhat} of \pycode{X} using the number of
PCs in part (ii).
\end{enumerate}

\begin{solution}
\begin{python}
# (i) Find PCs and mean
mu = np.mean(X, axis=0)
X_centered = X - mu
U, s, Vtr = svd(X_centered, full_matrices=False)
# PCs are the columns of Vtr (or rows of Vtr.T)
V = Vtr.T
# Eigenvalues are s^2 / (n-1)
eigenvals = s**2 / (len(X) - 1)

# (ii) Find minimum number of PCs for 90% variance
cumulative_var = np.cumsum(eigenvals) / np.sum(eigenvals)
nc = np.argmax(cumulative_var >= 0.90) + 1

# (iii) Create approximation
Z = X_centered @ V[:, :nc]
Xhat = Z @ V[:, :nc].T + mu
\end{python}
\end{solution}
\end{enumerate}
 
\end{document}

